\chapter{Open questions}
What this edition does not establish. These are recorded so a later pass knows where the evidence stopped.
\small
\begin{itemize}
\item What bit 7 of (0x60F020) means. Both GM handlers set it before posting their event, and nothing converted in this tree reads it.
\item What consumes the \{E=0x91, D=0x03, A=<(0x7F4D)>, W=0x04\} record on the 0x2E00 pending list. Until that is traced, 'GM System On replaces the current settings with the GM settings' is the panel dialog's claim, not something the disassembly establishes for the SysEx path.
\item What sets byte +1 of sub\_FB4CAE's argument record to 0x11 or 0x10 ? i.e. what actually makes the unit TRANSMIT GM System On/Off over MIDI (as opposed to writing it into an SMF). The template selection is proven; the trigger is not.
\item What (0x89) is and who sets it to 0xFF. While it is 0xFF, both SysEx drainers exit immediately and no SysEx at all is processed.
\item Whether trailing bytes between the sub-ID\#2 and the F7 are genuinely tolerated. Needs a hardware test: send F0 7E 7F 09 01 00 F7 and see whether GM mode turns on. Hardware is currently unreachable ? park it in kn7000\_mame/notes/HARDWARE-QUESTIONS-PENDING-FELIPE.md rather than asserting it.
\item Whether F0 50 7F 09 01 F7 really acts as GM System On (the ID byte is never re-checked after the front-end filter). Same hardware test applies.
\item Which company holds MIDI manufacturer ID 0x50. The tree refuses to name it without the MMA assignment list; the published manual should do the same.
\item What the Matsushita commands 1-33 do individually. Only the two GM slots were traced here; PtrTable\_F4F800 slot 9 (0xFB33FE) and slots 23-31 have distinct handlers that were not read.
\item Data\_F73844 and Data\_F75675 (prom\_b) are still classified as unidentified data blocks even though they are byte-identical SMF GM-event pairs with located readers. Their headers should be corrected in the disassembly; that is a source edit this pass did not make.
\item No reproducibility artefact was committed for this pass. The enumeration above was done by reading the .s files plus one throwaway Python byte-scan of original\_ROMs/*. If any of these byte sequences are published, a small committed probe (walk LinkTable\_F4FF61 from record 767 and print every accepted message) belongs beside the disassembly, per the project's permanent-storage rule.
\item The manual's 'BULK DUMP only on MIDI 1' is explained by call sites, not by a setting: sub\_FB7165 (send) names only ring 0x601432 and sub\_FB6098 (receive) names only ring 0x601646. But ORDINARY SysEx reception exists on both ports (sub\_FB21CB drains the MIDI 2 SysEx ring 0x601C6E every main-loop pass). If the manual also says non-bulk SysEx is MIDI-1-only, that is NOT what the firmware does and the
\item Which company holds manufacturer ID 0x50 is not settled in this tree, and the tree contradicts itself: the header at prom\_a 0xFA5AEB says 'Matsushita', notes/FINDINGS-midi-port.md explicitly refuses to name one. Needs the MMA assignment list.
\item That CPU 2's serial channel 0 is the physical MIDI 2 jack is an inference from the link-channel wiring (link selector 2 = CPU1->CPU2 MIDI bytes, link index 6 = CPU2->CPU1 MIDI bytes). The schematic pin table in notes/DRIVER-INSIGHT-wsa1-2026-09-02.md names only MIDI1OUT/MIDI1IN/MIDI1SNS, and that table has already needed one off-by-one correction. A MIDI2 pin label from the service manual would cl
\item sub\_FB63AF, which builds the 16-byte record that sub\_FB2160 dispatches on, was not traced, so the mapping from wire bytes to record fields 0 and 4 is unknown. Until that is done, 'command number = record[0], 0x00-0x21' cannot be turned into 'the byte at wire offset N is the command'. The 34-entry table at prom\_b 0xF4F800 is the next thing to enumerate for the manual.
\item Frame-length arithmetic is one byte from a contradiction and should be measured, not reasoned about: the transmitter stops adding nibble pairs once the frame reaches 0xFC = 252 bytes and then appends flag + checksum + F7, while the receiver aborts a message whose accumulated count reaches 0xFF = 255. Whether the largest frame the machine emits is 255 or 256 bytes - and therefore whether it can ove
\item The EXCLUSIVE field is edited with bit mask 0x0F but the transmit gates test only bit 3. The reading that ON writes all four bits (0x0F) and OFF writes 0x00 is consistent with both, but the bit-field setter at 0xFB... reached by `calr 0xf22b` from 0xF9AF37 was not read, so it is inference. Do not publish 'EXCLUSIVE = (0x7F38) bit 3' as the storage format without reading that setter.
\item sub\_FB7025 conditionally prepends the 3-byte fragment F0 50 7E when (0x60FCE4)+0x0F == 1, inside the same per-packet loop that already has a 12-byte header in the buffer. Whether that is the continuation-packet header (and therefore whether packets 2..n look different from packet 1 on the wire) is not established.
\item Reproducibility: the four-ROM address census quoted in the settings finding was run ad hoc as `python3` over original\_ROMs/*, scanning the 12 TLCS-900 direct-memory prefix spellings (0xC1/D1/E1/F1 with a 16-bit address field, 0xC2/D2/E2/F2 with a 24-bit one) for 0x7F32/33/35/36/38/39/3A/3B. notes/prom\_a\_xref.py and notes/prom\_c\_xrefs.py already do this properly for their own images; the cross-imag
\item Who reads (0xA9) bits 2 and 4? They are written (0x04 after a clean F7, 0x10 after an aborted SysEx) but a census of every (0xA9) site in prom\_a finds no test of either. If nothing reads them they are diagnostics only ? but a reader could live in prom\_c/prom\_d, which were not searched.
\item Do the command HANDLERS re-check the manufacturer byte? The grammar walker provably skips message bytes 0 and 1, so at the grammar level `F0 7E 21 04 00 11 F7` is accepted identically to `F0 50 21 ...`. Many handlers do re-read the raw buffer descriptor (0x60FC80), so a late check is possible; none was located. This must be settled before a manual states that either identifier works.
\item The tables contain command ids up to 0x27, but both dispatchers reject anything >= 0x22. Commands 0x22-0x27 therefore parse successfully and are then dropped by sub\_FB2160/sub\_FB225D. Is there a third dispatcher with a larger bound, or are those table entries dead? Only one site each references 0xF4F800 and 0xF4F888.
\item What do commands 0x01-0x1F actually do? On the interrupt-side ring, 0x01-0x08 and 0x0A-0x16 all land on the single shared handler 0xFB2820, which gates on the panel-mode cell (0x207A)==0x79 and on command >= 7; on the foreground ring the same commands land on a no-op. Distinct handlers exist only for 0x00, 0x09, 0x17, 0x18, 0x19, 0x1A, 0x1B-0x1F, 0x20, 0x21.
\item The 0x2B / 0x2C / 0x2D subtrees (338 reachable nodes) are the bulk of the protocol ? parameter addresses after `04 <00|01> 11` ? and were not enumerated. The committed probe dumps them with --paths; they need a pass of their own before any parameter table goes into a manual.
\item What is the unit byte after 0x04 for ? a device number selectable in the UI, or a MAIN/SUB distinction between two linked units? The grammar accepts 0x00 and 0x01 unconditionally and the value is only remembered (0x60FC91), never compared. Felipe's hardware could settle this by watching what the instrument sends for its own unit.
\item The 266-byte payload capacity per buffer is arithmetic on the 0x118 descriptor stride and the abutment at 0x60FC80, not a constant in the ROM. Worth confirming against a MAME memory map before it is printed as a limit; the byte-count bound of 255/256 bytes per message IS a ROM constant and is the number a manual should quote.
\item (0x0964) = 4, set by the foreground parser after a clean F7, has no located reader ? the same open question the existing block header at 0xFA5942 already flags.
\item The wording of ERROR 42 ("The data has not been received correctly") implies a handshake with the receiving MIDI device, but the only raise is a CPU1$\leftrightarrow$CPU2 link-block timeout (Link\_WaitBlockDone). The tree's own header for Link\_WaitBlockDone says of this very call site: "0xFB6FD1 ... what it waits on is NOT established". So whether the link block being waited on carries the outgoing MIDI stream to CP
\item Semantically odd but repeatedly confirmed: an unknown PRODUCT ID at message byte [2] gives status 0x07 $\rightarrow$ ERROR 41, while a wrong byte at [3], [4] or [5] gives ERROR 40. If a real-hardware check is ever possible, sending F0 50 99 04 00 11 ... (bogus product) vs F0 50 2D 99 00 11 ... (good product, bogus next byte) would distinguish 41 from 40 on the panel and settle it. Felipe's hardware is currently u
\item Trie depth 3 accepts 0x00 or 0x01 for every product id and both lead to the identical subtree. Reading this as a device/unit number is an INFERENCE (likely); the ROM only says the two values are accepted and equivalent.
\item The tail flag byte before the checksum must be 0x00 or 0x01 (status 0x12 otherwise) and is stored in field 15 of the parse record. What it MEANS is not established ? nothing converted reads field 15 back.
\item Statuses 0x19-0x1E each guard one bulk-dump section (their guard is `field 3 of *(0x60FCDC) == 0`, i.e. the previous section completed). Which section is which ? SOUND, COMBINATION, SYSTEM/PART \& MIDI, TOTAL KEYBOARD, per the menu text in Paint\_SysexBulkDump ? is NOT established and should not be guessed in the manual.
\item PtrTable\_F98DE5 (the German screen set) has no located reader; the tree's own header says so. Whether some other firmware variant of the WSA1R selects it, or it is dead weight from a shared source tree, is unknown.
\item Reproducibility: the new probe wsa1/notes/sysex-probes/sysex\_error\_codes.py + its README section are committed as kn5000-roms-disasm d4f11089 (with the required LLVM version line). It re-derives every number quoted above from original\_ROMs/, asserts both load bases and prints OK. The pre-existing wsa1/notes/sysex-probes/sysex\_grammar\_dump.py (0ac8f694) covers the trie. Nothing from this session li
\item Which incoming SysEx message actually selects a category for a remote-triggered dump. The five arms at 0xFB5122-0xFB514A are dispatch-table entries 0x1B..0x1F, but the receive trie at prom\_b 0xF5115B produces no form with those classes. Either a second grammar root exists or those classes are raised internally; neither was found. UNVERIFIED ? do not publish a "dump request" message format.
\item The class-6 arm of sub\_FB6026 (0xFB6041 `cp h,0x06` -> status 0x1F). No byte sequence in the trie yields class 6, so the wire form that triggers it is UNVERIFIED.
\item Press/release semantics. PanelButton\_Route (0xF861AC) passes the button index in W bits 0-4 and a flag in bit 7; the existing header reads bit 7 clear as "released", but that word is not cited to anything. On screen 0x79 the four row-select handlers act with bit 7 SET while the dump-start arm of 0xF99A8F acts with bit 7 CLEAR, and on the GENERAL MIDI screen (0xF99DDD, 0xF99DF4, 0xF99E2C) every han
\item What the constant bytes 0x04 and 0x11 mean. They are in every 21/22/2B/2C/2D message and the trie matches them literally. A product-family byte and a model id is the obvious reading but nothing in the ROM says so. Likewise nothing says whether a user can ever make the machine transmit device number 0x01 rather than 0x00; every transmit template holds 0x00 and only the receiver accepts either.
\item The address space. 0x080000 (SOUND), 0x100000 (SYSTEM,PART \& MIDI), 0x140000 (COMBINATION), 0x180000/0x188400 (SEQUENCER) are logical addresses that do not equal the internal source addresses (RAM 0x7600, remote flash 0xE80000/0xEC0000). Whether they map onto anything a user or a service manual would recognise is UNVERIFIED.
\item (0xC4). Its value 1 enables the SEQUENCER dump and 2 hides the SEQUENCER row entirely; a model or option variant is the obvious reading, but it is not established. This matters for the manual: on a machine where (0xC4) != 1 the SEQUENCER category is silently refused.
\item The second output ring 0x60153C, woken through T\_MIDI\_PostSendWork\_PortB (0xF40730 -> 0xFA5C12). It looks like a second MIDI OUT, and sub\_FB71BB sends to both rings, but no bulk-dump frame goes through it and the port's identity is UNVERIFIED.
\item The one-byte mailbox (0x9B), set to 0x00 by both sub\_FB7165 and sub\_FB71BB and to 0xF3 at 0xFE1507. notes/FINDINGS-midi-port.md records it as selecting an arm of MIDI\_PostSendWork but does not decode it.
\item The orphan at 0xFB248A and its template prom\_b 0xF4FF10 (F0 50 2D 04 00 11 20 00 00 00 00 10). Nothing names 0xFB248A, and its template says 16 bytes while the routine passes count 0x0000 to the link read. It should NOT go into the manual as a real message; the probe asserts the anomaly so it will fire if either half changes.
\item Frame timing on the wire. A 256-byte frame exactly fills the 0x100-byte MIDI-out ring and Ring\_Put\_0100 drops rather than blocks, so whether a real transfer ever loses a byte depends on interrupt latency, which cannot be settled from the disassembly. Felipe's hardware is unreachable (notes: kn5000-hardware-inaccessible), so this is a hardware question to park.
\end{itemize}
